\section{批判的な意見}
一方で，学習用データと実際のアプリケーションで必要とされるデータ分布を一致させるのは難しそうだと感じた．利用者が具体的に何をLLMに求め，日本語特有の知識がどこで必要になるのかを明確化できなければ，収集すべきコーパスの設計も曖昧になる．単に大量の日本語データを集めるだけでは不十分で，他言語では補えない固有知識をどう定義するかが重要になると思う．

また，オープンソース化は意義が大きい一方で，外部の企業が性能向上のためにデータを無制限に取り込むモデルを作り，結果として国内が培った資源が他者に大きく利用されるリスクもある．公的資金で運営する以上，日本の利用者に明確な利益が返る仕組みや，学習データの扱いに関する公平性を担保する枠組みが求められる．

さらに，現状ではデータ透明性が自主的な申告に依存しており，制度的に監査・開示を担保する仕組みは未成熟に見える．学習データの管理・公開を技術的に支援するシステムが整わない限り，倫理指針の運用や資金投入の基準作りは難しくなるはずだ．早期にその基盤が整備されることを期待したい．